\section{Related Work and Scope}
\label{sec:related-work}

EWC stores a quadratic approximation around earlier parameters and weights parameter motion by a Fisher-derived importance estimate \citep{EWC}. Other parameter-importance regularizers estimate importance from the optimization trajectory, output sensitivity, or a combination of EWC-style curvature and path-integral evidence \citep{SI,MAS,RWalk}; their apparently different estimators have also been analyzed within a common framework \citep{BenzingRegularisation}. Bayesian and posterior-combination approaches instead accumulate task posteriors or merge their moments \citep{ProgressCompress,VCL,IMM}. A practical ambiguity remains: implementations called ``Fisher'' may use the model-label expectation, sampled model labels, observed labels, or minibatch aggregates. These estimators are not interchangeable in theory or practice \citep{KunstnerEF,VanDeVenFisher2025}. Our scalar-loss derivation concerns the observed-label OPG object; sampled model-Fisher EWC is evaluated separately rather than treated as another name for the same matrix.

Regularization can instead protect functions or representations. LwF distills old predictions on later inputs, encoder-based lifelong learning constrains reconstructions of old-task features, and functional regularization preserves selected predictive functions or memorable points \citep{LwF,EBLL,FunctionalGP,FROMP}. More generally, many training objectives contain vectors or matrices that are computed from parameters or activations but held constant with respect to the parameters being optimized, such as weight-averaged teacher parameters \citep{MeanTeacher}, stop-gradient target projections \citep{BYOL,SimSiam}, predictions of an earlier model \citep{LwF}, and running activation moments \citep{TeReL}. EWC and IEWC follow the same pattern in the sense that both compute an importance matrix once earlier tasks are learned and then hold it fixed as the weighting of a quadratic penalty. CLOC complements this taxonomy: it begins with a preservation specification and asks which local stored penalty follows. The algebra below allows differentiable scalar functionals, while IEWC specializes to scalar-loss preservation and loss-level-set-normal output preservation.

iEF was introduced as a normalized observed-gradient approximation for natural-gradient optimization \citep{IEF}. We recover its Euclidean algebra from a different question: the metric distance of an old output displacement from the tangent to its old loss level set. This provides a preservation interpretation but does not convert iEF into exact global curvature. Our additions relative to simply substituting iEF into EWC are the CLOC derivation, the matched output-metric generalization, the finite-damping and interpolation boundary, and the separation of normalization from aggregation weights. The RDW ablation borrows a streaming residual-diversity intuition from GSS \citep{GSS}; it is an explicitly order-dependent heuristic, not part of the level-set result.
